\chapter{Our Data}
\label{ch:our-data}

\section{Data Sources}
\label{sec:data-sources}

Table~\ref{tab:data-sources} inventories every data source available for the project.
All sources span 11.3 years of history.

\begin{table}[htbp]
\centering
\caption{Available data sources, what they enable, and their constraints.}
\label{tab:data-sources}
\small
\begin{tabular}{%
  >{\raggedright\arraybackslash}p{2.8cm}
  >{\raggedright\arraybackslash}p{2.2cm}
  >{\raggedright\arraybackslash}p{4.0cm}
  >{\raggedright\arraybackslash}p{3.8cm}}
\toprule
\textbf{Data Source} & \textbf{Granularity} & \textbf{What It Enables} & \textbf{Key Constraint} \\
\midrule
Tick-level RV (34 symbols) &
  L1 tick &
  RV at any frequency, RQ, jumps, semivariances &
  --- \\[6pt]
Daily OHLCV (34 + VIX) &
  Daily &
  HAR baselines, ML training &
  No intraday structure \\[6pt]
E-mini L2 depth &
  L2 (${\sim}$4M ticks/day) &
  OBI, depth ratio, VPIN, LSTM input &
  E-mini only \\[6pt]
IV surfaces (Marquee ERDVOL) &
  Full tenor $\times$ strike grid &
  VRP, skew, term structure, butterfly &
  Full surface SPX only; single-stock ATM IV via EDRVOL\_PERCENT \\[6pt]
VIX term structure &
  Daily &
  Regime detection, contango/backwardation &
  --- \\[6pt]
Cross-asset (treasuries, FX, commodities) &
  Mixed (tick to daily) &
  Spillover features, macro regime &
  Daily for some instruments \\
\bottomrule
\end{tabular}
\end{table}


\section{GS Edge vs.\ Public Data}
\label{sec:gs-edge}

Table~\ref{tab:gs-edge} summarizes the capabilities our data provides relative to what is typically available in academic studies.

\begin{table}[htbp]
\centering
\caption{Data advantages over standard academic sources.}
\label{tab:gs-edge}
\small
\begin{tabular}{%
  >{\raggedright\arraybackslash}p{2.2cm}
  >{\raggedright\arraybackslash}p{3.0cm}
  >{\raggedright\arraybackslash}p{3.0cm}
  >{\raggedright\arraybackslash}p{3.8cm}}
\toprule
\textbf{Capability} & \textbf{Our Data} & \textbf{Public Alternative} & \textbf{Edge} \\
\midrule
RV estimation quality &
  Tick-level, 34 symbols, 11.3 years &
  5-min returns from TAQ / Oxford-Man &
  Precise RQ, jump detection, kernel estimators \\[6pt]
Options surface &
  Full SPX tenor $\times$ strike grid (Marquee) &
  VIX only (model-free 30-day) &
  Full surface derivatives: skew, butterfly, slope, event-implied \\[6pt]
Microstructure depth &
  E-mini L2 (4M ticks/day) &
  L1 quotes only (TAQ) &
  OBI at depth levels 2--5, true depth imbalance \\[6pt]
Cross-asset sync &
  Same tick timestamp across asset classes &
  Daily closes only &
  Intraday lead-lag detection \\[6pt]
Panel breadth &
  30 mega-caps + 4 ETFs + E-mini &
  Typically 1 index or 29 DJIA &
  Graph models, sector structure, cross-sectional features \\
\bottomrule
\end{tabular}
\end{table}

\begin{keyidea}[The data advantage is structural, not just bigger]
The tick-level RV, full IV surface, and L2 depth data let us construct features that are mathematically impossible to compute from public data. This is not a matter of having ``more data''; public researchers literally cannot replicate the feature set.
\end{keyidea}


\section{Constraints That Shape Decisions}
\label{sec:data-constraints}

L2 depth data covers the E-mini only, so microstructure depth features (OBI at levels 2--5, depth ratio) apply exclusively to the index.
Equities receive L1 features only.
The full IV surface is SPX only, so surface-derived features (skew, term structure, butterfly) function as market-wide regime signals.
Single-stock ATM IV is available via Marquee EDRVOL\_PERCENT, enabling per-name VRP and IV--RV gap features.
These constraints determine which features apply to which assets in the feature engineering pipeline.

\begin{warning}[Do not treat depth or surface features as stock-level predictors]
L2 depth and the full IV surface are index-level data. Using them as if they vary across individual equities would introduce look-ahead bias or a constant signal masquerading as stock-specific information.
Single-stock ATM IV is the exception---it is genuinely per-name.
\end{warning}
